\section{\S 3.3 电解池}\label{sect:电解池}

    \subsection{电解基本原理}\label{subsec:电解基本原理}

        \textcolor{red}{电解：将电能转化为化学能的装置（使非自发的氧化还原反应得以实现）。}

        \textbf{阴极}：与外电源负极相连——发生还原反应

        \textbf{阳极}：与外电源正极相连——发生氧化反应

        电解的三个过程：

        \begin{enumerate}
            \item \textbf{电离}：电解质产生自由移动的正负离子
            \item \textbf{迁移}：正负离子分别移向阴、阳两极
            \item \textbf{放电}：正负离子分别在两极上得失电子
        \end{enumerate}

        \begin{xmp}[电解 CuCl2]{电解 \cemh{CuCl2} 溶液}
            阴极（\cemh{Cu^{2+}} 得电子）：\cemh{Cu^{2+} + 2e- -> Cu}

            阳极（\cemh{Cl-} 失电子）：\cemh{2Cl- - 2e- -> Cl2 ^}

            总反应：\cemh{CuCl2} \xleq[电解]{} \cemh{Cu + Cl2 ^}
        \end{xmp}

    \subsection{离子的放电顺序}\label{subsec:离子放电顺序}

        \subsubsection{阴极放电顺序（阳离子得电子）}\label{subsubsec:阴极放电顺序}

            \textcolor{red}{\cemh{Ag+} > \cemh{Fe^{3+}} > \cemh{Cu^{2+}} > \cemh{H+}}（酸）\textcolor{red}{> \cemh{Fe^{2+}} > \cemh{Zn^{2+}} > \cemh{H+}}（水）\textcolor{red}{> \cemh{Al^{3+}} > \cemh{Mg^{2+}} > \cemh{Na+} > \cemh{Ca^{2+}} > \cemh{K+}}

            \textcolor{blue}{注：\cemh{H+} 在金属活动性顺序中氢前铝后看浓度。}

        \subsubsection{阳极放电顺序（阴离子失电子）}\label{subsubsec:阳极放电顺序}

            惰性电极（\cemh{C}、\cemh{Pt}、\cemh{Au}）：

            \textcolor{red}{活性电极 > \cemh{I-} > \cemh{Br-} > \cemh{Cl-} > \cemh{OH-} > (族价)含氧酸根 > \cemh{F-}}

            \textcolor{blue}{活性电极（除 \cemh{Pt}、\cemh{Au} 外的金属）优先于阴离子放电——阳极溶解。}

        \begin{xmp}[电解 Na2SO4]{电解 \cemh{Na2SO4} 溶液——本质电解水}
            阴极：\cemh{2H2O + 2e- -> H2 ^ + 2OH-}（\cemh{H+} > \cemh{Na+}）

            阳极：\cemh{2H2O - 4e- -> O2 ^ + 4H+}（\cemh{OH-} > \cemh{SO4^{2-}}）

            总反应：\cemh{2H2O} \xleq[电解]{} \cemh{2H2 ^ + O2 ^}
        \end{xmp}

        \begin{xmp}[电解 NaCl]{电解饱和食盐水}
            阴极：\cemh{2H2O + 2e- -> 2OH- + H2 ^}

            阳极：\cemh{2Cl- - 2e- -> Cl2 ^}

            总反应：\cemh{2NaCl + 2H2O} \xleq[电解]{} \cemh{2NaOH + H2 ^ + Cl2 ^}

            \textcolor{blue}{鲁科版：强调离子放电。沪科版：强调与总反应关系。}
        \end{xmp}

        \begin{xmp}[电解 CuSO4]{电解 \cemh{CuSO4} 溶液}
            阴极：\cemh{Cu^{2+} + 2e- -> Cu}

            阳极：\cemh{2H2O - 4e- -> O2 ^ + 4H+}

            总反应：\cemh{2CuSO4 + 2H2O} \xleq[电解]{} \cemh{2Cu + O2 ^ + 2H2SO4}
        \end{xmp}

    \subsection{电极对电极反应的影响}\label{subsec:电极对电极反应的影响}

        \begin{xmp}[活性阳极]{\cemh{Cu} 作阳极}
            以 \cemh{Cu} 作阳极电解 \cemh{H2SO4} 溶液：

            阳极区：\cemh{Cu - 2e- -> Cu^{2+}}（溶液变蓝色）

            阴极区：\cemh{2H+ + 2e- -> H2 ^}（产生气泡）

            总反应：\cemh{Cu + 2H+} \xleq[通电]{} \cemh{Cu^{2+} + H2 ^}

            \textcolor{red}{电解使非自发进行的氧化还原反应得以实现。}

            \textcolor{blue}{若阴极把 \cemh{Cu} 换成石墨：阴极区先产生气泡，后析出红色固体（\cemh{Cu^{2+}} 与 \cemh{H+} 竞争的结果）。}
        \end{xmp}

    \subsection{电解的应用}\label{subsec:电解的应用}

        \subsubsection{精炼铜}\label{subsubsec:精炼铜}

            \textbf{阳极}：粗铜（\cemh{Fe}、\cemh{Zn}、\cemh{Ni} 等活泼金属溶解，\cemh{Ag}、\cemh{Au} 形成阳极泥）

            \textbf{阴极}：精铜（\cemh{Cu^{2+} + 2e- -> Cu} 沉积）

            \textbf{电解质}：含 \cemh{Cu^{2+}} 的溶液

            \begin{itemize}
                \item 阳极：\cemh{Zn - 2e- -> Zn^{2+}}、\cemh{Fe - 2e- -> Fe^{2+}}、\cemh{Ni - 2e- -> Ni^{2+}}、\cemh{Cu - 2e- -> Cu^{2+}}
                \item 阴极：\cemh{Cu^{2+} + 2e- -> Cu}
            \end{itemize}

            \textcolor{blue}{\cemh{Cu^{2+}} 浓度略降（以\SI{2}{\mole} \cemh{e-} 为例，阳极溶解少于阴极沉积）。}

        \subsubsection{电镀}\label{subsubsec:电镀}

            \textbf{阳极}：镀层金属

            \textbf{阴极}：镀件

            \textbf{电镀液}：含镀层金属离子的溶液

            \textcolor{red}{电镀液的浓度不变，镀层更均匀。}

            \begin{xmp}[电镀银]{铜上镀银}
                阳极：\cemh{Ag}（镀层金属）

                阴极：\cemh{Cu}（镀件）

                电镀液：\cemh{AgNO3} 溶液

                阳极：\cemh{Ag - e- -> Ag+}

                阴极：\cemh{Ag+ + e- -> Ag}

                电镀液浓度：\textcolor{red}{不变}
            \end{xmp}

        \subsubsection{活性阳极的其他应用}\label{subsubsec:活性阳极其他应用}

            \begin{xmp}[废水处理]{含 \cemh{Cr2O7^{2-}} 酸性废水处理}
                以 \cemh{Fe} 为阳极电解处理含 \cemh{Cr2O7^{2-}} 的酸性废水。

                \textbf{阳极}（\cemh{Fe}）：\cemh{Fe - 2e- -> Fe^{2+}}（\cemh{Fe^{2+}} 现制现用）

                \textbf{阴极}：\cemh{2H+ + 2e- -> H2 ^}

                溶液中：\cemh{6Fe^{2+} + Cr2O7^{2-} + 14H+ -> 6Fe^{3+} + 2Cr^{3+} + 7H2O}

                然后增大溶液 pH，\cemh{Cr(OH)3 v} + \cemh{Fe(OH)3 v} 沉淀除去。

                \textcolor{blue}{加入 \cemh{NaCl} 增强溶液导电性（\cemh{Cl-} 在阳极不放电，因为 \cemh{Fe} 是活性电极优先溶解）。}
            \end{xmp}

            \textcolor{blue}{阳极氧化：通过电解，在阳极给金属（通常是 \cemh{Al}）镀上氧化膜的工艺。}

        \subsubsection{氯碱工业}\label{subsubsec:氯碱工业}

            电解饱和食盐水：

            \begin{center}
            \begin{tabular}{ll}
                阳极： & \cemh{2Cl- - 2e- -> Cl2 ^} \\
                阴极： & \cemh{2H2O + 2e- -> 2OH- + H2 ^} \\
                总反应： & \cemh{2NaCl + 2H2O} \xleq[]{电解} \cemh{2NaOH + H2 ^ + Cl2 ^}
            \end{tabular}
            \end{center}

            \textcolor{red}{使用阳离子交换膜，防止 \cemh{Cl2} 与 \cemh{NaOH} 反应。}

        \subsubsection{海水淡化}\label{subsubsec:海水淡化}

            电解结合离子交换膜技术（电渗析），实现海水淡化。

        \subsubsection{电化学制备}\label{subsubsec:电化学制备}

            \begin{xmp}[CO2 资源化]{电解 \cemh{CO2} 制 \cemh{HCOOH}}
                \textbf{酸性溶液}：

                阴极：\cemh{CO2 + 2H+ + 2e- -> HCOOH}

                阳极：\cemh{2H2O - 4e- -> 4H+ + O2 ^}

                \textbf{\cemh{HCO3-} 溶液}：

                阴极：\cemh{CO2 + HCO3- + 2e- -> HCOO- + CO3^{2-}}

                阳极：\cemh{4HCO3- - 4e- -> 4CO2 ^ + O2 ^ + 2H2O}
            \end{xmp}

    \subsection{电解过程中的 pH 变化}\label{subsec:电解pH变化}

        \textcolor{red}{极区 pH 变化看电极反应，不是离子迁移。}

        \textcolor{red}{极区 pH 变化 \textcolor{black}{$\neq$} 溶液 pH 变化。}

        \begin{center}
        \small
        \begin{tabular}{|c|c|c|c|}
            \hline
            \textbf{电解液} & \textbf{阴极区} & \textbf{阳极区} & \textbf{溶液 pH}\\
            \hline
            \cemh{NaCl} & pH 增大（生成 \cemh{OH-}） & pH 减小 & 增大 \\
            \hline
            \cemh{CuSO4} & pH 不变 & pH 减小（生成 \cemh{H+}） & 减小 \\
            \hline
            \cemh{Na2SO4} & pH 增大 & pH 减小 & 不变 \\
            \hline
            \cemh{HCl} & pH 增大（消耗 \cemh{H+}） & pH 不变 & 增大 \\
            \hline
        \end{tabular}
        \end{center}
